\begin{vidu}%[3P7-6.2B][Loai1]
Cho phản ứng nhiệt hạch: $\Hn{1}{2}{D}+\Hn{1}{2}{D}\to\Hn{2}{4}{He}$ tỏa năng lượng 23,7 MeV. Biết độ hụt khối của hạt nhân \Hn{1}{2}{D} là 0,0025u. Lấy $1u=931,5\ MeV/c^2$. Năng lượng liên kết của hạt nhân \Hn{2}{4}{He} bằng
\choice
{21,3 MeV}
{26,0 MeV}
{\True 28,4 MeV}
{19,0 MeV}
\loigiai{
Năng lượng phản ứng tỏa ra
\begin{align*}
W = \left(\Delta m_{He} - 2\Delta m_D \right)c^2 \Rightarrow W_{lkHe} = W + 2\Delta m_D.c^2 = 23,7 + 2.0,0025.931,5 = 28,4\ MeV.
\end{align*}
}
\end{vidu}
\begin{vidu}%[3P7-6.2B][Loai1]
Biết phản ứng nhiệt hạch: $\Hn{1}{2}{D}+\Hn{1}{2}{D}\to \Hn{2}{3}{He}+n$ tỏa ra một năng lượng 3,25 MeV. Độ hụt khối của $\Hn{1}{2}{D}$ là 0,0024u. Biết $1uc^2=931,5\ MeV$. Năng lượng liên kết của hạt nhân $\Hn{2}{3}{He}$ là
\choice
{5,22 MeV}
{9,24 MeV}
{8,52 MeV}
{\True 7,72 MeV}
\loigiai{
Năng lượng tỏa ra là:
\begin{align*}
&W=\left(\Delta m_{He}-2\Delta {m_D}\right). c^2=3,25\ MeV \Rightarrow \Delta m_{He}=8,289.10^{-3}u\\
\Rightarrow &{W_{lkHe}}=\Delta m_{He}. c^2=8,289.10^{-3}. 931,5=7,72\ MeV.
\end{align*}	
}
\end{vidu}
\begin{vidu}%[3P7-6.2K][Loai1]
Cho phản ứng hạt nhân: $\Hn{1}{2}{D}+\Hn{1}{2}{D}\to \Hn{2}{3}{He}+\Hn{0}{1}{n}$. Biết khối lượng của $\Hn{1}{2}{D}$,  $\Hn{2}{3}{He}$, $\Hn{0}{1}{n}$ lần lượt là: ${m_D}=2,0135u;\ m_{He}=3,0149u;\ {m_n}=1,0087u$. Cho khối lượng hạt nhân bằng số khối; $N_A=6,02.10^{23}\ mol^{-1}$. Lấy $1u=931,5\ MeV/c^2$. Khối lượng Deuterium cần thiết để có thể thu được năng lượng nhiệt hạch tương đương với năng lượng toả ra khi đốt 1 tấn than là (biết năng lượng toả ra khi đốt 1 kg than là 30000 kJ)
\choice
{\True 0,4 g}
{4 kg}
{8 g}
{4 g}
\loigiai{
\begin{itemize}
\item Năng lượng tỏa ra trong một phản ứng là: $W= \left[2{m_D}-(m_{He}+{m_n}) \right]c^2=3,1671\ MeV$.
\item Năng lượng toả ra khi phản ứng nhiệt hạch hết m gam \Hn{1}{2}{D} (tương đương với N hạt D): 
\begin{align*}
Q_1=\dfrac{W}{2}N=\dfrac{W}{2}.\dfrac{m}{M}.N_A.
\end{align*}
\item Tổng năng lượng tỏa ra khi đốt 1 tấn than là: $Q_2=10^3.30000.10^3=3.10^{10}J=1,875.10^{23}\ MeV$.
\item Theo bài ra: $Q_1=Q_2 \Leftrightarrow \dfrac{W}{2}.\dfrac{m}{M}.N_A=Q_2 \Rightarrow m=\dfrac{2MQ_2}{W.N_A}=0,4\ g$.
\end{itemize}
}
\end{vidu}
\begin{vidu}%[3P7-6.2G][Loai1]
Cho phản ứng nhiệt hạch: $\Hn{1}{2}{D}+\Hn{1}{2}{D}\to \Hn{2}{3}{He}+\Hn{0}{1}{n}$, biết độ hụt khối của $\Hn{1}{2}{D}$ và $\Hn{2}{3}{He}$ lần lượt là 0,0024u và 0,0305u. Nước trong tự nhiên có khối lượng riêng là: $1000\ kg/m^3$ và lẫn 0,015\% $D_2O$. Cho khối lượng hạt nhân bằng số khối; $N_A=6,02.10^{23}\ mol^{-1}$. Lấy $1u=931,5\ MeV/c^2$. Nếu toàn bộ $\Hn{1}{2}{D}$ được tách ra từ $1\ m^3$ nước trong tự nhiên làm nhiên liệu cho phản ứng trên thì năng lượng tỏa ra là
\motcot
{$1,89.10^{26}\ MeV$}
{\True $1,08.10^{26}\ MeV$}
{$1,02.10^{26}\ MeV$}
{$1,96.10^{26}\ J$}
\loigiai{
\begin{itemize}
\item Năng lượng tỏa ra trong một phản ứng là: $W=\left( \Delta m_{He}-2\Delta {m_D} \right)c^2=23,94\ MeV$.
\item Khối lượng $D_2O$ có trong $1\ m^3$ nước là: 
\begin{align*}
&{{m}_{{D_2}O}}=0,015\%m_{H_2O}=0,015\%. D_{H_2O}. V_{H_2O}=0,15\ kg=150\ g\\
\Rightarrow\ &N_{D_2O}=n_{D_2O}. {N_A}=\dfrac{m_{D_2O}}{M_{D_2O}}.N_A=\dfrac{150}{16+2.2}. 6,02.10^{23}=4,515.10^{24}\ \ct{hạt}.
\end{align*}
\item Số hạt Deuterium có $1\ m^3$ nước là: ${N_D}=2N_{D_2O}=9,03.10^{24}$ hạt.
\item Tổng năng lượng tỏa ra là: $Q=\dfrac{N_D}{2}W=1,0812.10^{26}\ MeV$.
\end{itemize}
}
\end{vidu}
\begin{vidu}%[3P7-6.2G][Loai1]
\nguon{QG - 2016} Giả sử ở một ngôi sao, sau khi chuyển hóa toàn bộ hạt nhân hydrogen thành hạt nhân $\Hn{2}{4}{He}$ thì ngôi sao lúc này chỉ có $\Hn{2}{4}{He}$ với khối lượng $4,6.10^{32}$ kg. Tiếp theo đó, $\Hn{2}{4}{He}$ chuyển hóa thành hạt nhân $\Hn{6}{12}{C}$ thông qua quá trình tổng hợp: $\Hn{2}{4}{He}+\Hn{2}{4}{He}+\Hn{2}{4}{He}\to \Hn{6}{12}{C}+7,27 \ MeV$. Coi toàn bộ năng lượng tỏa ra từ quá trình tổng hợp này đều được phát ra với công suất trung bình là $5,3.10^{30}$ W. Cho biết: 1 năm bằng 365,25 ngày, khối lượng mol của $\Hn{2}{4}{He}$ là 4 g/mol, số Avogadro $N_A = 6,02.10^{23}\ mol^{-1}$, $1eV=1,6.10^{-19}\ J$. Thời gian để chuyển hóa hết $\Hn{2}{4}{He}$ ở ngôi sao này thành $\Hn{6}{12}{C}$ vào khoảng
\motcot
{481,5 triệu năm}
{481,5 nghìn năm}
{160,5 nghìn năm}
{\True 160,5 triệu năm}
\loigiai{
\begin{itemize}
\item Số hạt nhân He trong $m = 4,6.10^{32}$ kg là: $N=\dfrac{m}{M}.N_A=\dfrac{4,6.10^{32}.10^3}{4}.6,02.10^{23}=6,923.10^{58}$.
\item Cứ 1 phản ứng cần 3 hạt nhân He nên số phản ứng cho đến khi hết He là: $N_0=\dfrac{N}{3} $.
\item Năng lượng tỏa ra cho đến khi hết helium là $Q=N_0.7,27.1,6.10^{-13}=\dfrac{N}{3}.7,27.1,6.10^{-13}\ J $.
\item Thời gian để chuyển hóa hết helium là 
\begin{align*}
t&=\dfrac{Q}{P}=\dfrac{N.7,27.1,6.10^{-13}}{3P}=\dfrac{6,923.10^{58}.7,27.1,6.10^{-13}}{3.5,3.10^{30}}\\
&=5,065.10^{15}\ s=1,605.10^8\ \ct{năm} = 160,5\ \ct{triệu năm}.
\end{align*}
\end{itemize}
}
\end{vidu}

